\clearpage\fancyhead[L]{\small LEMNISCATISCHE FUNCTIONEN II.}\pdfbookmark[0]{De curva lemniscata}{de-curva-lemniscata}\SourceLeaf{422}

\begin{center}
\textbf{DE CURVA LEMNISCATA.}\par\medskip
\rule{2cm}{0.3pt}
\end{center}

\CurvaArticle{1}{\textbf{1.}}
Posito integrali $\int\frac{\mathrm{d}x}{\surd(1-x^4)}$, a $x=0$ usque ad $x=s$, $=\varphi$, dicimus $s$ sinum lemniscaticum ipsius $\varphi$, $s=\sin\operatorname{lemn}\varphi$.

\CurvaArticle{2}{\textbf{2.}}
Valor integralis ab $x=0$ usque ad $x=1$ est $=1.3110287771\;4605987$ secundum \textsc{Stirling}, qui valor a nobis usque ad figuram undecimam verus in ventus est, utentibus formula: $\operatorname{arc}\sin\operatorname{lemn}\frac{7}{23}+2\operatorname{arc}\sin\operatorname{lemn}\frac12$ (\textsc{Euler} habet $1.311031$). Potestates huius numeri, cuius duplum semper per $\varpi$ designabimus, has invenimus
\[
\begin{array}{r@{\quad}c@{\quad}r@{\quad}l@{\quad}l}
1& \mathord{.}\mathord{.}\mathord{.}\mathord{.} &1.3110287771&4605990680&320.7\\[5pt]
2& \mathord{.}\mathord{.}\mathord{.}\mathord{.} &1.7187964545&0509311.7&\\[5pt]
4& \mathord{.}\mathord{.}\mathord{.}\mathord{.} &2.9542612520&1927863.4&\\[5pt]
5& \mathord{.}\mathord{.}\mathord{.}\mathord{.} &3.8731215170&0712625\;4&\\[5pt]
6& \mathord{.}\mathord{.}\mathord{.}\mathord{.} &5.0777737656&5251025.3&\\[5pt]
8& \mathord{.}\mathord{.}\mathord{.}\mathord{.} &8.7276595451&8251569.0&\\[5pt]
9& \mathord{.}\mathord{.}\mathord{.}\mathord{.} &11.4422128208&59&\\[5pt]
12& \mathord{.}\mathord{.}\mathord{.}\mathord{.} &25.7837864151&41749&\\[5pt]
13& \mathord{.}\mathord{.}\mathord{.}\mathord{.} &33.8032859402&5&
\end{array}
\]


\SourceLeaf{423}
\[
\begin{aligned}
\operatorname{log}\operatorname{brigg}\dfrac12\varpi&=0.1176122226\quad9692.2\\[5pt]
\operatorname{log}\operatorname{hyp}\dfrac12\varpi&=0.2708121550\quad7159155410\quad6425\\[5pt]
\operatorname{log}\operatorname{hyp}\varpi&=0.9639593356\quad3153686352\quad36577
\end{aligned}
\]
Sinum lemniscaticum ipsius $(\tfrac12\varpi-a)$ cosinum lemniscaticum ipsius $a$ dicemus.

\CurvaArticle{3}{\textbf{3.}}
Aequationis
\[
\dfrac{\mathrm{d}x}{\surd(1-x^4)}+\dfrac{\mathrm{d}y}{\surd(1-y^4)}=0
\]
integrale completum invenitur hoc
\[
\dfrac{x\surd\dfrac{1-yy}{1+yy}+y\surd\dfrac{1-xx}{1+xx}}
{1-xy\surd\dfrac{1-xx}{1+xx}\,\dfrac{1-yy}{1+yy}}=C
\]

Sed eiusdem aequat. integrale est
\[
\operatorname{arc}\sin\operatorname{lemn}x+\operatorname{arc}\sin\operatorname{lemn}y=c
\]
unde sequitur $C$ esse functionem ipsius $c$. Ut appareat qualis, ponamus $y=0$, tum fit $C=x$, $c=\operatorname{arc}\sin\operatorname{lemn}x$, quare erit $c=\operatorname{arc}\sin\operatorname{lemn}C$, sive $C=\sin\operatorname{lemn}c$. Hinc si $\sin\operatorname{lemn}p=x$, $\sin\operatorname{lemn}q=y$, erit
\[
\sin\operatorname{lemn}(p+q)=
\dfrac{x\surd\dfrac{1-yy}{1+yy}+y\surd\dfrac{1-xx}{1+xx}}
{1-xy\surd\dfrac{1-xx}{1+xx}\,\dfrac{1-yy}{1+yy}}
\]

Hinc posito $p=\tfrac12\varpi$, $q=-a$, fit propter
\[
\sin\operatorname{lemn}(-a)=-\sin\operatorname{lemn}a,\quad
\sin\operatorname{lemn}\dfrac12\varpi=1
\]
\[
\cos\operatorname{lemn}a=\surd\dfrac{1-\sin\operatorname{lemn}a^2}{1+\sin\operatorname{lemn}a^2}
\]
Forma autem praecedens transit in hanc


\SourceLeaf{424}
\[
\begin{aligned}\sin\operatorname{lemn}(p\pm q)&=\\[7pt]&\quad{}
\dfrac{\sin\operatorname{lemn}p\,\cos\operatorname{lemn}q\pm\sin\operatorname{lemn}q\,\cos\operatorname{lemn}p}
{1\mp\sin\operatorname{lemn}p\,\sin\operatorname{lemn}q\,\cos\operatorname{lemn}p\,\cos\operatorname{lemn}q}\end{aligned}
\]
\[
\begin{aligned}\cos\operatorname{lemn}(p\pm q)&=\\[7pt]&\quad{}
\dfrac{\cos\operatorname{lemn}p\,\cos\operatorname{lemn}q\mp\sin\operatorname{lemn}p\,\sin\operatorname{lemn}q}
{1\pm\sin\operatorname{lemn}p\,\sin\operatorname{lemn}q\,\cos\operatorname{lemn}p\,\cos\operatorname{lemn}q}\end{aligned}
\]

\noindent[Spätere Bemerkung:]

\noindent I.\hfill $\alpha+\beta+\gamma=180^0\;[=\varpi]$\hfill\mbox{}

Setzt
\[
\begin{array}{ll}
\sin\operatorname{lemn}\alpha=\operatorname{tang}a,&\cos\operatorname{lemn}\alpha=\cos A\\[5pt]
\sin\operatorname{lemn}\beta=\operatorname{tang}b,&\cos\operatorname{lemn}\beta=\cos B\\[5pt]
\sin\operatorname{lemn}\gamma=\operatorname{tang}c,&\cos\operatorname{lemn}\gamma=\cos C
\end{array}
\]
so sind $a,\ b,\ c,\ A,\ B,\ C$ Seiten und Winkel eines sphärischen Dreiecks, wo
\[
\dfrac{\sin A}{\sin a}=\dfrac{\sin B}{\sin b}=\dfrac{\sin C}{\sin c}=\surd2
\]

\par\Needspace{7\baselineskip}\noindent II.\hfill $\alpha+\beta+\gamma=90^0\;[=\tfrac12\varpi]$\hfill\mbox{}

Setzt
\[
\begin{array}{ll}
\sin\operatorname{lemn}\alpha=\cos a,&\cos\operatorname{lemn}\alpha=-\operatorname{tang}A\\[5pt]
\sin\operatorname{lemn}\beta=\cos b,&\cos\operatorname{lemn}\beta=-\operatorname{tang}B\\[5pt]
\sin\operatorname{lemn}\gamma=\cos c,&\cos\operatorname{lemn}\gamma=-\operatorname{tang}C
\end{array}
\]
so sind wieder $a,\ b,\ c,\ A,\ B,\ C$ Seiten und Winkel eines sphärischen Dreiecks, in welchem
\[
\dfrac{\sin A}{\sin a}=\dfrac{\sin B}{\sin b}=\dfrac{\sin C}{\sin c}=\surd\dfrac12
\]

\CurvaArticle{4}{[4.]}
Si valores $s\;[=\sin\frac{\pi}{\varpi}\varphi]$, qui reddunt ipsum $\sin\operatorname{lemn}\varphi=0$ secundum regulas notas productum infinitum generare concipiuntur, nec non valores ipsius $s$, qui reddunt ipsum $\sin\operatorname{lemn}\varphi=\infty$, quorum primum sit $P\varphi$, secundum $Q\varphi$, permissum erit (id quod rigorose demonstrare possumus) ponere
\[
\sin\operatorname{lemn}\varphi=\dfrac{P\varphi}{Q\varphi}
\]


\SourceLeaf{425}
\noindent erit vero
\[
\begin{aligned}P\varphi &=\alpha s
 \Bigl(1+\dfrac{4ss}{(e^{\pi}-e^{-\pi})^{2}}\Bigr)
 \Bigl(1+\dfrac{4ss}{(e^{2\pi}-e^{-2\pi})^{2}}\Bigr)
 \\[9pt]&\quad{}\Bigl(1+\dfrac{4ss}{(e^{3\pi}-e^{-3\pi})^{2}}\Bigr)\mathord{.}\mathord{.}\mathord{.}\end{aligned}
\]
\[
\begin{aligned}Q\varphi &=
 \Bigl(1-\dfrac{4ss}{(e^{\frac12\pi}+e^{-\frac12\pi})^{2}}\Bigr)
 \Bigl(1-\dfrac{4ss}{(e^{\frac32\pi}+e^{-\frac32\pi})^{2}}\Bigr)
 \\[9pt]&\quad{}\Bigl(1-\dfrac{4ss}{(e^{\frac52\pi}+e^{-\frac52\pi})^{2}}\Bigr)\mathord{.}\mathord{.}\mathord{.}\end{aligned}
\]
\[
\begin{aligned}\alpha &=\dfrac{\varpi}{\pi}\end{aligned}
\]
Simili modo positis
\[
\begin{aligned}p\varphi &=c
 \Bigl(1-\dfrac{4ss}{(e^{\pi}+e^{-\pi})^{2}}\Bigr)
 \Bigl(1-\dfrac{4ss}{(e^{2\pi}+e^{-2\pi})^{2}}\Bigr)
 \\[9pt]&\quad{}\Bigl(1-\dfrac{4ss}{(e^{3\pi}+e^{-3\pi})^{2}}\Bigr)\mathord{.}\mathord{.}\mathord{.}\end{aligned}
\]
\[
\begin{aligned}q\varphi &=
 \Bigl(1+\dfrac{4ss}{(e^{\frac12\pi}-e^{-\frac12\pi})^{2}}\Bigr)
 \Bigl(1+\dfrac{4ss}{(e^{\frac32\pi}-e^{-\frac32\pi})^{2}}\Bigr)
 \\[9pt]&\quad{}\Bigl(1+\dfrac{4ss'}{(e^{\frac52\pi}-e^{-\frac52\pi})^{2}}\Bigr)\mathord{.}\mathord{.}\mathord{.}\end{aligned}
\]
\noindent erit
\[
\cos\operatorname{lemn}\varphi=\dfrac{p\varphi}{q\varphi}
\]
\CurvaArticle{5}{[5.]}
\[
\begin{aligned}
{}^{4}\!\surd 2\, .\,P(\varphi+\dfrac12\varpi)&=p\varphi\\[5pt]
{}^{4}\!\surd 2\, .\,Q(\varphi+\dfrac12\varpi)&=q\varphi
\end{aligned}
\]
\[
\begin{aligned}
p(\varphi+\dfrac12\varpi)&=-{}^{4}\!\surd 2\, .\,P\varphi\\[5pt]
q(\varphi+\dfrac12\varpi)&=\phantom{-}{}^{4}\!\surd 2\, .\,Q\varphi
\end{aligned}
\]
\[
\begin{aligned}
Pi\psi\varpi&=ie^{\pi\psi\psi}P\psi\varpi\\[5pt]
Qi\psi\varpi&=e^{\pi\psi\psi}Q\psi\varpi\\[5pt]
pi\psi\varpi&=e^{\pi\psi\psi}q\psi\varpi\\[5pt]
qi\psi\varpi&=e^{\pi\psi\psi}p\psi\varpi
\end{aligned}
\]
\noindent [werden die an einem andern Orte untersuchten Reihen Seite 405 d.\ B.
\[
\begin{aligned}
&\varphi-\dfrac1{60}\varphi^{5}-\dfrac1{10080}\varphi^{9}+\mathord{.}\mathord{.}\\[5pt]
&1+\dfrac1{12}\varphi^{4}-\dfrac1{10080}\varphi^{8}+\mathord{.}\mathord{.}\\[5pt]
&1-\dfrac12\varphi\varphi-\dfrac1{24}\varphi^{4}-\mathord{.}\mathord{.}\mathord{.}\\[5pt]
&1+\dfrac12\varphi\varphi-\dfrac1{24}\varphi^{4}+\mathord{.}\mathord{.}\mathord{.}
\end{aligned}
\]
\noindent resp.\ mit $\mathfrak{P}\varphi$, $\mathfrak{Q}\varphi$, $\mathfrak{p}\varphi$, $\mathfrak{q}\varphi$ bezeichnet, so ist:]
\[
\begin{aligned}
\mathfrak{P}\psi\varpi&=e^{\frac12\pi\psi\psi}P\psi\varpi
\\[5pt] \mathfrak{p}\psi\varpi&=e^{\frac12\pi\psi\psi}p\psi\varpi\\[5pt]
\mathfrak{Q}\psi\varpi&=e^{\frac12\pi\psi\psi}Q\psi\varpi
\\[5pt] \mathfrak{q}\psi\varpi&=e^{\frac12\pi\psi\psi}q\psi\varpi
\end{aligned}
\]



\SourceLeaf{426}
\CurvaArticle{6}{[6.]}
Ex expressionibus supra allatis sequitur
\[
\begin{aligned}
\sin\operatorname{lemn}\varphi
&=\dfrac{\dfrac{\pi}{\varpi}(e^{\frac12\pi}-e^{-\frac12\pi})}
        {e^{\frac12\pi}-2s+e^{-\frac12\pi}}
 -\dfrac{\dfrac{\pi}{\varpi}(e^{\frac12\pi}-e^{-\frac12\pi})}
        {e^{\frac12\pi}+2s+e^{-\frac12\pi}}\\[5pt]
&\quad-\dfrac{\dfrac{\pi}{\varpi}(e^{\frac32\pi}-e^{-\frac32\pi})}
        {e^{\frac32\pi}-2s+e^{-\frac32\pi}}
 +\dfrac{\dfrac{\pi}{\varpi}(e^{\frac32\pi}-e^{-\frac32\pi})}
        {e^{\frac32\pi}+2s+e^{-\frac32\pi}}\\[5pt]
&\quad+\text{etc.}\\[5pt]
\cos\operatorname{lemn}\varphi&= 
\end{aligned}
\]
Hinc vero sequitur
\[
\begin{aligned}\sin\operatorname{lemn}\psi\varpi
&=\dfrac{\pi}{\varpi}\dfrac{4}{e^{\frac12\pi}+e^{-\frac12\pi}}\sin\psi\pi
\\[7pt]&\quad{}-\dfrac{\pi}{\varpi}\dfrac{4}{e^{\frac32\pi}+e^{-\frac32\pi}}\sin3\psi\pi
+\mathord{.}\mathord{.}\mathord{.}\end{aligned}
\]
\CurvaArticle{7}{[7.]}
\[
\begin{aligned}\log(1+\mu\cos\varphi)
&=2\Bigl\{\dfrac{\mu}{1+\surd(1-\mu\mu)}\cos\varphi
\\[7pt]&\quad{}-\Bigl(\dfrac{\mu}{1+\surd(1-\mu\mu)}\Bigr)^{2}\dfrac{\cos2\varphi}{2}
\\[7pt]&\quad{}+\Bigl(\dfrac{\mu}{1+\surd(1-\mu\mu)}\Bigr)^{3}\dfrac{\cos3\varphi}{3}\end{aligned}
\]
\[
\begin{aligned}
\log Q\psi\varpi&=-\dfrac12\log2+\dfrac1{12}\pi\\[5pt]
&\quad+\dfrac{2}{e^{\pi}-e^{-\pi}}\cos2\psi\pi
-\dfrac12\dfrac{2}{e^{2\pi}-e^{-2\pi}}\cos4\psi\pi
\\[7pt]&\quad{}+\dfrac13\dfrac{2}{e^{3\pi}-e^{-3\pi}}\cos6\psi\pi-\mathord{.}\mathord{.}
\end{aligned}
\]
\[
\begin{aligned}
\log P\psi\varpi&=\dfrac12\log2-\dfrac16\pi+\log\sin\psi\pi\\[5pt]
&\quad-\dfrac{2}{e^{2\pi}-1}\cos2\psi\pi
-\dfrac12\dfrac{2}{e^{4\pi}-1}\cos4\psi\pi
\\[7pt]&\quad{}-\dfrac13\dfrac{2}{e^{6\pi}-1}\cos6\psi\pi-\mathord{.}\mathord{.}
\end{aligned}
\]
\[
\begin{aligned}
\log q\psi\varpi&=-\dfrac14\log2+\dfrac1{12}\pi\\[5pt]
&\quad-\dfrac{2}{e^{\pi}-e^{-\pi}}\cos2\psi\pi
-\dfrac12\dfrac{2}{e^{2\pi}-e^{-2\pi}}\cos4\psi\pi
\\[7pt]&\quad{}-\dfrac13\dfrac{2}{e^{3\pi}-e^{-3\pi}}\cos6\psi\pi-\mathord{.}\mathord{.}
\end{aligned}
\]
\[
\begin{aligned}
\log p\psi\varpi&=\dfrac34\log2-\dfrac16\pi+\log\cos\psi\pi\\[5pt]
&\quad+\dfrac{2}{e^{2\pi}-1}\cos2\psi\pi
-\dfrac12\dfrac{2}{e^{4\pi}-1}\cos4\psi\pi
\\[7pt]&\quad{}+\dfrac13\dfrac{2}{e^{6\pi}-1}\cos6\psi\pi-\mathord{.}\mathord{.}
\end{aligned}
\]
\[
\begin{aligned}
\log\sin\operatorname{lemn}\psi\varpi&=\log2-\dfrac14\pi+\log\sin\psi\pi\\[5pt]
&\quad-\dfrac{2}{e^{\pi}-1}\cos2\psi\pi
+\dfrac12\dfrac{2}{e^{2\pi}-1}\cos4\psi\pi
\\[7pt]&\quad{}-\dfrac13\dfrac{2}{e^{3\pi}-1}\cos6\psi\pi-\mathord{.}\mathord{.}
\end{aligned}
\]



\SourceLeaf{427}
\[
\begin{array}{r@{\;}c@{\;}r@{\;}l@{\quad}l}
\log Q\psi\pi={}&-&0.0847742372&7&\\[5pt]
&+&0.0865895371&57&\cos2\psi\pi\\[5pt]
&-&0.0018674144&52&\cos4\psi\pi\\[5pt]
&+&0.0000537996&86&\cos6\psi\pi\\[5pt]
&-&17436&71&\cos8\psi\pi\\[5pt]
&+&602&81&\cos10\psi\pi\\[5pt]
&-&21&71&\cos12\psi\pi\\[5pt]
&+&&81&\cos14\psi\pi\\[5pt]
&-&&\phantom{0}3&\cos16\psi\pi
\end{array}
\]
\[
\begin{array}{r@{\;}c@{\;}r@{\;}l@{\quad}l}
\log P\psi\varpi={}&\multicolumn{4}{l}{\log\sin\psi\pi}\\[5pt]
&-&0.1770251853&2&\\[5pt]
&-&0.0037418731&98&\cos2\psi\pi\\[5pt]
&-&34873&54&\cos4\psi\pi\\[5pt]
&-&43&42&\cos6\psi\pi\\[5pt]
&-&&\phantom{0}6&\cos8\psi\pi
\end{array}
\]
\CurvaArticle{8}{[8.]}
\[
\begin{aligned}\sin\operatorname{lemn}\psi\varpi
&=\\[8pt]&\quad{}\surd\dfrac{4}{e^{\frac12\pi}}\, .\,
\dfrac{\sin\psi\pi-e^{-2\pi}\sin3\psi\pi+e^{-6\pi}\sin5\psi\pi-\mathord{.}\mathord{.}}
{1+2e^{-\pi}\cos2\psi\pi+2e^{-4\pi}\cos4\psi\pi-\mathord{.}\mathord{.}}\end{aligned}
\]
\[
\begin{aligned}P\psi\varpi&=2^{\frac34}\surd\dfrac{\pi}{\varpi}\, .\,
\\[7pt]&\quad{}\bigl\{e^{-\frac14\pi}\sin\psi\pi-e^{-\frac94\pi}\sin3\psi\pi
+e^{-\frac{25}{4}\pi}\sin5\psi\pi-\mathord{.}\mathord{.}\bigr\}\end{aligned}
\]
\[
\begin{aligned}Q\psi\varpi&=\dfrac{1}{2^{\frac14}}\surd\dfrac{\pi}{\varpi}\, .\,
\\[7pt]&\quad{}\bigl\{1+2e^{-\pi}\cos2\psi\pi+2e^{-4\pi}\cos4\psi\pi-\mathord{.}\mathord{.}\bigr\}\end{aligned}
\]
\[
\begin{aligned}
1-2e^{-\pi}+2e^{-4\pi}-2e^{-9\pi}+\mathord{.}\mathord{.}\mathord{.}\mathord{.}
&=\surd\dfrac{\varpi}{\pi}\\[5pt]
e^{-\frac14\pi}+e^{-\frac94\pi}+e^{-\frac{25}{4}\pi}+\mathord{.}\mathord{.}\mathord{.}\mathord{.}\mathord{.}\mathord{.}
&=\dfrac12\surd\dfrac{\varpi}{\pi}
\end{aligned}
\]
\[
\begin{aligned}
\surd\dfrac{\varpi}{\pi}&=0.9135791381\;5611682140\;724259\\[5pt]
2e^{-\frac14\pi}&=0.9118762555\;3199247353\;1842589456\;058838833\\[5pt]
2e^{-\frac94\pi}&=0.0017028766\;8561031607\;1704906\\[5pt]
2e^{-\frac{25}{4}\pi}&=.\,.\,.\,.\,.\,.\,.\,.\,.59\;3851399312\;9644497731\;18\\[5pt]
2e^{-\frac{49}{4}\pi}&=.\,.\,.\,.\,.\,.\,.\,.\,.\,.\,.\;.\,.\,.\,.\,.\,.3867\;40505991
\end{aligned}
\]



\SourceLeaf{428}
\[
\begin{aligned}
2e^{-\pi}&=0.0864278365\ 2754449954\ 8835474343\ 4560225514\\[5pt]
2e^{-4\pi}&=0.0000069746\ 8471241799\ 0983550387\ 96535\\[5pt]
2e^{-9\pi}&=\text{. .......... .105109703 5201288}\\[5pt]
2e^{-16\pi}&=\text{. .......... .......... .295807}
\end{aligned}
\]
\CurvaArticle{9}{[9.]}
\[
\begin{aligned}\dfrac{1}{\sin\operatorname{lemn}\psi\varpi}
&=\dfrac{\pi}{\varpi}\biggl(\dfrac{1}{\sin\psi\pi}
-\dfrac{4}{e^{\pi}+1}\sin\psi\pi
\\[7pt]&\quad{}-\dfrac{4}{e^{3\pi}+1}\sin3\psi\pi
-\dfrac{4}{e^{5\pi}+1}\sin5\psi\pi-.\,.\biggr).\end{aligned}
\]
\[
\begin{aligned}
\dfrac{1}{P\psi\varpi}
&=\dfrac{\pi}{\varpi}\biggl[\dfrac{1}{\sin\psi\pi}
\\[7pt]&\quad{}-4(e^{-2\pi}-e^{-6\pi}+e^{-12\pi}-\mathord{.}\mathord{.}\mathord{.})\sin\psi\pi\\[5pt]
&-4(e^{-4\pi}-e^{-10\pi}+e^{-18\pi}-\mathord{.}\mathord{.}\mathord{.})\sin3\psi\pi\\[5pt]
&-4(e^{-6\pi}-e^{-14\pi}+e^{-24\pi}-\mathord{.}\mathord{.}\mathord{.})\sin5\psi\pi\\[5pt]
&-.\,.
\end{aligned}
\]
\CurvaArticle{10}{[10].}
\[
\begin{aligned}
\sin\operatorname{lemn}\psi\varpi&=\operatorname{tang}u\pi\\[5pt]
\pi\,\mathrm{d}u&=\varpi\cos\operatorname{lemn}\psi\varpi.\mathrm{d}\psi
\end{aligned}
\]
\[
\begin{aligned}\cos\operatorname{lemn}\psi\varpi
&=\dfrac{\pi}{\varpi}\biggl(
\dfrac{4}{e^{\frac12\pi}+e^{-\frac12\pi}}\cos\psi\pi
\\[7pt]&\quad{}+\dfrac{4}{e^{\frac32\pi}+e^{-\frac32\pi}}\cos3\psi\pi+\mathord{.}\mathord{.}\mathord{.}\biggr)\end{aligned}
\]
\[
\begin{aligned}u\pi&=\dfrac{4}{e^{\frac12\pi}+e^{-\frac12\pi}}\sin\psi\pi
\\[7pt]&\quad{}+\dfrac13\dfrac{4}{e^{\frac32\pi}+e^{-\frac32\pi}}\sin3\psi\pi+\mathord{.}\mathord{.}\mathord{.}\end{aligned}
\]
\[
1+2p\cos\varphi+2pp\cos2\varphi+\mathord{.}\mathord{.}\mathord{.}
=\dfrac{\dfrac1p-p}{\dfrac1p+p-2\cos\varphi}
\]
\[
1-2p\cos\varphi+2pp\cos2\varphi-\mathord{.}\mathord{.}\mathord{.}
=\dfrac{\dfrac1p-p}{\dfrac1p+p+2\cos\varphi}
\]
\[
p\cos\varphi+p^3\cos3\varphi+\mathord{.}\mathord{.}\mathord{.}
=\dfrac{(\dfrac1p-p)\cos\varphi}{(\dfrac1p+p)^2-4\cos\varphi^2}
\]
\[
2p\sin\varphi+\dfrac23p^3\sin3\varphi+\mathord{.}\mathord{.}\mathord{.}
=\operatorname{arc\,tang}\dfrac{2\sin\varphi}{\dfrac1p-p}
\]
\[
\begin{aligned}\dfrac12u\pi
&=\operatorname{arc\,tang}\dfrac{2\sin\psi\pi}{e^{\frac12\pi}-e^{-\frac12\pi}}
\\[7pt]&\quad{}-\operatorname{arc\,tang}\dfrac{2\sin\psi\pi}{e^{\frac32\pi}-e^{-\frac32\pi}}+\mathord{.}\mathord{.}\mathord{.}\end{aligned}
\]




\SourceLeaf{429}
\CurvaArticle{11}{[11.]}
\[
\begin{aligned}
(P\psi\varpi)^2
&=\dfrac{\surd\dfrac\pi\varpi}{2^{\frac74}\cos\dfrac18\pi}
\\[6pt]&\quad{}(1+2e^{-2\pi}\cos4\psi\pi+2e^{-8\pi}\cos8\psi\pi+\mathord{.}\mathord{.}\mathord{.})\\[4pt]
&\quad-\dfrac{\surd\dfrac\pi\varpi}{2^{\frac74}\sin\dfrac18\pi}
\\[7pt]&\quad{}(2e^{-\frac12\pi}\cos2\psi\pi+2e^{-\frac92\pi}\cos6\psi\pi
+2e^{-\frac{25}2\pi}\cos10\psi\pi+\mathord{.}\mathord{.}\mathord{.})
\end{aligned}
\]
\[
\begin{aligned}
(Q\psi\varpi)^2
&=2^{\frac34}\cos\dfrac18\pi.\surd\dfrac\pi\varpi.
(1+2e^{-2\pi}\cos4\psi\pi+\mathord{.}\mathord{.}\mathord{.}\\[5pt]
&\quad+2^{\frac34}\sin\dfrac18\pi.\surd\dfrac\pi\varpi.
(2e^{-\frac12\pi}\cos2\psi\pi+\mathord{.}\mathord{.}\mathord{.}
\end{aligned}
\]
\[
A=\dfrac{\surd\dfrac\pi\varpi}{2^{\frac72}\cos\dfrac18\pi},
\qquad
B=\dfrac{\surd\dfrac\pi\varpi}{2^{\frac74}\sin\dfrac18\pi}
\]
\[
\begin{aligned}
0.3522376226\ 6118372314&=A\\[5pt]
0.3535519576\ 3585935635&=2Be^{-\frac12\pi}\\[5pt]
0.0013155679\ 2352259042&=2Ae^{-2\pi}\\[5pt]
0.0000012329\ 5741446398&=2Be^{-\frac92\pi}\\[5pt]
\text{. .......... 0856752170}&=2Ae^{-8\pi}\\[5pt]
\text{. .......... ......1494}&=2Be^{-\frac{25}2\pi}
\end{aligned}
\]
\CurvaArticle{12}{[12.]}
\begin{center}Variae Summationes serierum absconditae.\end{center}
\[
\begin{aligned}1^{0}\quad&
\biggl[\dfrac{2}{e^\pi+e^{-\pi}}\biggr]^2
+\biggl[\dfrac{2}{e^{2\pi}+e^{-2\pi}}\biggr]^2
+\biggl[\dfrac{2}{e^{3\pi}+e^{-3\pi}}\biggr]^2+\mathord{.}\mathord{.}\mathord{.}
\\[8pt]&=\dfrac12\dfrac{\varpi\varpi}{\pi\pi}-\dfrac{1}{2\pi}-\dfrac12\end{aligned}
\]
\[
\begin{aligned}2^{0}\quad&
\biggl[\dfrac{2}{e^{\frac12\pi}-e^{-\frac12\pi}}\biggr]^2
+\biggl[\dfrac{2}{e^{\frac32\pi}-e^{-\frac32\pi}}\biggr]^2
+\biggl[\dfrac{2}{e^{\frac52\pi}-e^{-\frac52\pi}}\biggr]^2+\mathord{.}\mathord{.}\mathord{.}
\\[8pt]&=\dfrac12\dfrac{\varpi\varpi}{\pi\pi}-\dfrac{1}{2\pi}\end{aligned}
\]
\[
\begin{aligned}3^{0}\quad&
\biggl[\dfrac{2}{e^\pi-e^{-\pi}}\biggr]^2
+\biggl[\dfrac{2}{e^{2\pi}-e^{-2\pi}}\biggr]^2
+\biggl[\dfrac{2}{e^{3\pi}-e^{-3\pi}}\biggr]^2+\mathord{.}\mathord{.}\mathord{.}
\\[8pt]&=\dfrac16-\dfrac{1}{2\pi}\end{aligned}
\]
\[
\begin{aligned}4^{0}\quad&
\biggl[\dfrac{2}{e^{\frac12\pi}+e^{-\frac12\pi}}\biggr]^2
+\biggl[\dfrac{2}{e^{\frac32\pi}+e^{-\frac32\pi}}\biggr]^2
+\biggl[\dfrac{2}{e^{\frac52\pi}+e^{-\frac52\pi}}\biggr]^2+\mathord{.}\mathord{.}\mathord{.}
\\[8pt]&=\dfrac{1}{2\pi}\end{aligned}
\]
\[
\begin{aligned}5^{0}\quad&
\dfrac{2}{e^{\frac12\pi}+e^{-\frac12\pi}}
+\dfrac{2}{e^{\frac32\pi}+e^{-\frac32\pi}}+\mathord{.}\mathord{.}\mathord{.}
\\[8pt]&=\dfrac12\dfrac\varpi\pi\end{aligned}
\]
\nopagebreak[4]
\[
\begin{aligned}6^{0}\quad&\end{aligned}
\]



\SourceLeaf{430}
\CurvaArticle{13}{[13.]}
\begin{center}\textit{Rechnungen zur Lemniscata gehörig.}\end{center}
\[
\begin{aligned}
\{(\sin\operatorname{lemn}\dfrac25\varpi)^2
+(\sin\operatorname{lemn}\dfrac45\varpi)^2\}^2&=14\surd5-30\\[5pt]
\{(\sin\operatorname{lemn}\dfrac25\varpi)^2
-(\sin\operatorname{lemn}\dfrac45\varpi)^2\}^2&=10\surd5-22
\end{aligned}
\]
\[
\begin{gathered}
\dfrac12(\sin\operatorname{lemn}\dfrac25\varpi)^4
+\dfrac12(\sin\operatorname{lemn}\dfrac45\varpi)^4
\\[5pt]=6\surd5-13\\[5pt]
=0.4164078649\ 9873817845\ 5042012387\ 65741
\end{gathered}
\]
\[
\begin{gathered}
\{\dfrac12(\sin\operatorname{lemn}\dfrac25\varpi)^4
-\dfrac12(\sin\operatorname{lemn}\dfrac45\varpi)^4\}^4
\\[5pt]=340-152\surd5\\[5pt]
=0.1176674200\ 3196614580\ 5602352846\ 01228
\end{gathered}
\]
\noindent daraus Radix
\[
0.3430268503\ 0761971797\ 7310507555\ 85731
\]
\noindent also
\[
\begin{aligned}(\sin\operatorname{lemn}\dfrac45\varpi)^4
&=\\[5pt]&\quad{}0.0733810146\ 9111846047\ 7731504831\ 80010\end{aligned}
\]
\[
\begin{aligned}(\sin\operatorname{lemn}\dfrac25\varpi)^4
&=\\[5pt]&\quad{}0.7594347153\ 0635789643\ 2352519943\ 51472\end{aligned}
\]
\[
\begin{aligned}(\sin\operatorname{lemn}\dfrac45\varpi)^2
&=\\[5pt]&\quad{}0.2708893033\ 8999814497\ 30710\end{aligned}
\]
\[
\begin{aligned}(\sin\operatorname{lemn}\dfrac25\varpi)^2
&=\\[5pt]&\quad{}0.8714555153\ 9155336074\ 646029\end{aligned}
\]
\[
\begin{aligned}
Q\dfrac1{10}\varpi
&={}^{5}\!\surd(\dfrac34\cos\dfrac1{10}\pi+\dfrac14\sin\dfrac1{10}\pi)\\[5pt]
Q\dfrac2{10}\varpi
&={}^{5}\!\surd(\dfrac34\cos\dfrac1{10}\pi-\dfrac14\sin\dfrac1{10}\pi)
\end{aligned}
\]
\CurvaArticle{14}{[14.]}
\begin{center}\textit{Lemniscatische Function.}\end{center}
\[
\int\dfrac{\mathrm{d}x}{\surd(1-x^4)}=\varphi,
\qquad
\int\dfrac{\mathrm{d}X}{\surd(1-X^4)}=\Phi
\]
\[
\begin{array}{rcccc}
\text{Es verhalte sich}&1&x&\surd(1-xx)&\surd(1+xx)\\[5pt]
\text{wie}&p&q&r&s\\[7pt]
\text{und}&1&X&\surd(1-XX)&\surd(1+XX)\\[5pt]
\text{wie}&P&Q&R&S\quad.
\end{array}
\]
\noindent so verhalten sich die $\varphi+\Phi$ entsprechenden Grössen


\SourceLeaf{431}
\begingroup\small\setlength{\arraycolsep}{4pt}\[
\begin{array}{@{}ll|l|l|l@{}}
\text{wie}&ppPP+qqQQ&pqRS+rsPQ&prPR-qsQS&psPS+qrQR\\
\text{oder}&pqRS-rsPQ&qqPP-ppQQ&qrPS-psQR&qsPR-prQS\\
\text{oder}&prPR+qsQS&qrPS+psQR&ppRR-qqSS&2pqPQ+rsRS\\
&&&=rrPP-ssQQ&\\
\text{oder}&psPS-qrQR&qsRR+prQS&rsRS-2pqPQ&ssPP+rrQQ\\
&&&&=ppSS+qqRR
\end{array}
\]\endgroup
\noindent Es sei
\[
\begin{aligned}
\sin\operatorname{lemn}X&=x\\[5pt]
\sin\operatorname{lemn}Y&=y\\[5pt]
\sin\operatorname{lemn}Z&=z
\end{aligned}
\]
\noindent und $X,\ Y,\ Z$ so von einander abhängig, dass
\[
X+Y+Z=0
\]
\noindent Man setze
\[
\int xx\,\mathrm{d}X=FX
\]
\noindent und
\[
FX+FY+FZ=u
\]
\noindent Aus früher vorgekommenen Formeln folgt
\[
\begin{aligned}
xx-zz&=yz\surd(1-x^4)-xy\surd(1-z^4)\\[5pt]
yy-zz&=xz\surd(1-y^4)-xy\surd(1-z^4)\\[5pt]
xx-yy&=yz\surd(1-x^4)-xz\surd(1-y^4)
\end{aligned}
\]
\noindent Aus
\[
\begin{aligned}
\mathrm{d}u&=\dfrac{xx\,\mathrm{d}x}{\surd(1-x^4)}
+\dfrac{yy\,\mathrm{d}y}{\surd(1-y^4)}
+\dfrac{zz\,\mathrm{d}z}{\surd(1-z^4)}\\[5pt]
0&=\dfrac{\mathrm{d}x}{\surd(1-x^4)}
+\dfrac{\mathrm{d}y}{\surd(1-y^4)}
+\dfrac{\mathrm{d}z}{\surd(1-z^4)}
\end{aligned}
\]
\noindent folgt
\[
\begin{aligned}
\mathrm{d}u&=(yy-xx)\dfrac{\mathrm{d}y}{\surd(1-y^4)}
+(zz-xx)\dfrac{\mathrm{d}z}{\surd(1-z^4)}\\[5pt]
&=xz\,\mathrm{d}y+xy\,\mathrm{d}z-yz.\surd(1-x^4).
\\[7pt]&\quad{}\biggl[\dfrac{\mathrm{d}y}{\surd(1-y^4)}+\dfrac{\mathrm{d}z}{\surd(1-z^4)}\biggr]\\[5pt]
&=xz\,\mathrm{d}y+xy\,\mathrm{d}z+yz\,\mathrm{d}x=\mathrm{d}.xyz
\end{aligned}
\]
